\documentclass[12pt]{report}
\usepackage[utf8]{inputenc}
\usepackage[T1]{fontenc}
\usepackage[a4paper, left=2.4cm, right=2.4cm, top=2cm, bottom=2cm, nomarginpar]{geometry}
\usepackage{amsmath, algorithm, pgfplots, pgf, pdfpages, amsfonts, amssymb, stmaryrd, enumitem,graphicx, algpseudocode, setspace, listings, courier, color, caption, centernot, hyperref, lstautogobble, zi4, mathtools, lmodern, colortbl, array, multicol, fancyhdr, lastpage, svg}
\usepackage{booktabs}
\hypersetup{
    pdftitle={CE204_FP1_G6},
    pdftoolbar=true,        % show Acrobat’s toolbar?
    pdfmenubar=true,        % show Acrobat’s menu?
    pdfauthor={Muhammet Yağcıoğlu},
    pdfsubject={CE204_FP_RPs},
    pdfcreator={Muhammet Yağcıoğlu},   % creator of the document
    pdfproducer={Muhammet Yağcıoğlu}, % producer of the document
    colorlinks,
    linkcolor={black},
    citecolor={black},
    urlcolor={blue!80!black}
}
\usepackage[noblocks]{authblk}
\usepackage[version=4]{mhchem}
\usepackage[export]{adjustbox}
\usepackage{eso-pic} % Required for absolute positioning
\graphicspath{ {./images/} }
\usetikzlibrary{patterns}
\usepackage{tikz-3dplot}
\usepackage[version=4]{mhchem}
\usepackage[export]{adjustbox}
\usepgfplotslibrary{groupplots}
\usepackage{xcolor}
\usepackage[absolute,overlay]{textpos}
\usepackage{ragged2e}
\usepackage{xcolor}
\usepackage{amsthm}
\usepackage{tikz}




%----------compiler settings---------
\setlength{\parindent}{0pt}
\everymath{\displaystyle}
\onehalfspacing

%-----------colors----------
\pgfplotsset{compat=1.18}
\definecolor{bluekeywords}{rgb}{0.13, 0.13, 1}
\definecolor{greencomments}{rgb}{0, 0.5, 0}
\definecolor{redstrings}{rgb}{0.9, 0, 0}
\definecolor{graynumbers}{rgb}{0.5, 0.5, 0.5}
\definecolor{bronze}{rgb}{0.82, 0.41, 0.12}
\definecolor{crimson}{rgb}{0.86, 0.08, 0.24}
\definecolor{brown(web)}{rgb}{0.65, 0.16, 0.16}
\definecolor{darkspringgreen}{rgb}{0.16, 0.32, 0.75}
\definecolor{blue(ncs)}{rgb}{0.0, 0.53, 0.74}
\definecolor{amber}{rgb}{1.0, 0.49, 0.0}
\definecolor{cadmiumgreen}{rgb}{0.0, 0.42, 0.24}




%-------code frame-----------
\lstset{
    autogobble,
    columns=fullflexible,
    showspaces=false,
    showtabs=false,
    breaklines=true,
    showstringspaces=false,
    breakatwhitespace=true,
    escapeinside={(*@}{@*)},
    commentstyle=\color{greencomments},
    keywordstyle=\color{bluekeywords},
    stringstyle=\color{redstrings},
    numberstyle=\color{graynumbers},
    basicstyle=\ttfamily\footnotesize,
    frame=l,
    framesep=12pt,
    xleftmargin=12pt,
    tabsize=4,
    captionpos=b
}



\lstdefinelanguage{myMMA}{
keywords={SetDirectory, NotebookDirectory, Exp, IdentityMatrix, Eigenvalues, 
ListPlot, PlotRange, PlotStyle, Directive, PointSize, AspectRatio, Blue, Graphics, Line, 
Manipulate, Show, Sqrt, UniformDistribution, GammaDistribution, BetaDistribution, 
Nintegrate, For, DataRange, AxesLabel, PlotLabel, Transpose, Export, Plot, Append, Infinity},
keywordstyle=\color{black},
commentstyle=\color{gray}, 
identifierstyle=\color{blue},
sensitive=false,
comment=[l]{(*},
morecomment=[s]{/*}{*/},
morestring=[b]',
morestring=[b]"
}






%-----gauss declare---------
\pgfmathdeclarefunction{gauss}{2}{%
  \pgfmathparse{1/(#2*sqrt(2*pi)) * exp(-((x-#1)^2)/(2*#2^2))}%
}
\usetikzlibrary{arrows.meta, decorations.markings}
\pgfplotsset{compat=1.17}







%--------------page numbering------------
\fancyfoot[R]{page \thepage\ of \textcolor{bronze}{\pageref{LastPage}}}
\fancyhead[R]{Group 6}



%-------------qframe settings---------------

\usepackage[framemethod=TikZ]{mdframed}
\newmdenv[%
    skipabove=1em, % space above the frame
    skipbelow=1em, % space below the frame
    linewidth=1pt, % width of the frame lines
    linecolor=gray!80, % color of the frame lines
    backgroundcolor=gray!9, % background color inside the frame
    roundcorner=5pt, % radius of the rounded corners
    innerleftmargin=1em, % margin within the frame at the left
    innerrightmargin=1em, % margin within the frame at the right
    innertopmargin=0.5em, % margin within the frame at the top
    innerbottommargin=0.5em % margin within the frame at the bottom
]{qframe}

\newenvironment{q}
{
    \begin{qframe}
    \noindent\textit{\textbf{Problem Statement:}}
    \par\smallskip
}
{
    \end{qframe}
}





%-------------ANSWER TAG---------------
\newcommand{\AnswerTag}{\hfill \begin{tikzpicture} \fill[black] (0.35cm,0) -- (0,0.175cm) -- (0.35cm,0.35cm) -- cycle;\end{tikzpicture}}



%-------Theorem and another question frames----------
\newenvironment{theorem}{\begin{mdframed}\textbf{Theorem.}}{\end{mdframed}}

%usage: \begin{theorem}

\newenvironment{question}{\begin{mdframed}\textbf{Question.}}{\end{mdframed}}

%usage: %usage: \begin{question}
\fancyhead[L]{CE204 Engineering Surveying - Report of Field Practice 3}
\author{Muhammet Yağcıoğlu - 290204042}
\usepackage{hyperref} 

\begin{document}
\onecolumn
\thispagestyle{empty}
\begin{center}


\textbf{\fontsize{14}{\baselineskip}\selectfont İzmir Institute of Technology}


\textbf{\fontsize{14}{\baselineskip}\selectfont In collaboration with}


\textbf{\fontsize{14}{\baselineskip}\selectfont Department of Civil Engineering}

\bigskip

\includesvg[width=5cm]{Picture1.svg}

\bigskip

\textbf{\fontsize{15}{\baselineskip}\selectfont CE204 – ENGINEERING
SURVEYING}

\bigskip

\textbf{\fontsize{16}{\baselineskip}\selectfont 2023-2024 SPRING}

\bigskip

\bigskip



\textit{\textbf{\fontsize{14}{\baselineskip}\selectfont Field Practice – 3}}

\bigskip



\bigskip

\textbf{\fontsize{14}{\baselineskip}\selectfont Lecturer: Assoc. Prof. Volkan Emre UZ }

\bigskip

\textbf{\fontsize{14}{\baselineskip}\selectfont Group 6}

\bigskip

\textbf{\fontsize{14}{\baselineskip}\selectfont Team Leader: Melih Berke	KARAADAM}

\bigskip

Melih Berke	KARAADAM\\
Mert BAY\\
Mert BOSTANCI \\
Mevlüt Yağız KÜL\\
Muhammet YAĞCIOĞLU\\
Muhammet Taylan ARSLAN\\


\bigskip

\bigskip

\bigskip

{\today}
\end{center}

\clearpage



\newpage
\thispagestyle{fancy}
\pagestyle{fancy}
\begin{abstract}
This paper elaborates on the methods and results of a thorough field practice carried out during the spring of 2023–2024 as part of the Engineering Surveying (CE204) course. Two crucial exercises, the Two-Peg Test and the Stadia Method for determining the stadia constant (\(K\)), were the focus of our research, which sought to master and validate the accuracy of levelling devices. To verify that our levelling device was accurate and free of systematic errors, we used the Two-Peg Test. This allowed us to take more trustworthy elevation difference readings. Simultaneously, our levelling equipment was further shown to be versatile and precise by allowing the empirical measurement of horizontal distances using the Stadia Method. Our research shows that conventional surveying methods are effective in improving the quality of topographical surveys, and that accurate instruments are crucial. By engaging in this practice, we improved our technical skills and made a substantial contribution to engineering surveying knowledge. This practice demonstrates the practical and educational benefits of doing rigorous fieldwork in civil engineering.
\end{abstract}
\tableofcontents
\newpage


\section*{Introduction}
\addcontentsline{toc}{section}{Introduction}

As part of an extensive fieldwork activity, this paper details the procedures followed and the results of accurate topographical surveying jobs. The primary goals were to determine the level's stadia constant (\(K\)) by using the stadia technique for distance measuring and to assess the levelling instrument's accuracy using the two-peg test. 

\subsection*{Part I: Two-Peg Test for Level Accuracy Verification}
For precise height determination, the reliability of levelling devices is of utmost importance. As a critical assessment approach, the two-peg test is based on the idea that, from each perspective, the vertical error margin should be the same when the instrument is positioned equidistantly between two fixed rods. By following this technique, any systematic faults in the instrument's sighting capabilities may be easily detected and corrected, leading to accurate elevation difference readings.

\subsection*{Part II: Stadia Method for Determining the Stadia Constant (\(K\))}
When using a level to measure horizontal distances empirically, the stadia technique is used. To transform the differential reading of the top and bottom stadia wires into an exact horizontal distance, this method requires computing the stadia constant (\(K\)) of a particular instrument. This technique demonstrates how levelling instruments can be used for more than only measuring elevation; they can also be used to measure distances directly across the measured area.

Improving the precision and dependability of topographical surveys relies heavily on both parts of the fieldwork, which in turn helps students of civil engineering learn more about measuring procedures, equipment calibration, and error correction.

\section*{Procedure}
\addcontentsline{toc}{section}{Procedure}

The fieldwork approach was created to solve particular goals related to the accuracy and utility of levelling devices used in topographical surveying. It was split into two phases.

\subsection*{Part I: Verification of Leveling Instrument Accuracy via the Two-Peg Test}

\begin{enumerate}

\item  \textit{Site Preparation and Measurement.} A reasonably level piece of ground was used to establish two points, A and B, 40 metres apart. A steel tape was used to confirm the exact distance.

\item \textit{Instrument Placement.} Precisely in the centre (C) between points A and B was where the levelling device was placed. By positioning themselves so that observations were equally far from both sites, parallax errors were reduced to a minimum.

\item \textit{Leveling Accuracy.} The device was carefully levelled such that the spherical bubble remained in the exact centre of its vial for the whole 360-degree revolution. Eliminating any levelling mistakes during measuring was dependent on this step.

\item \textit{Data Collection.} Accurate to within 1 mm, the staff recorded the values at positions A and B. A plumb bob was used to check the staff's verticality in order to prevent any systematic inaccuracies.

\item \textit{Secondary Positioning and Measurement.} After that, it was necessary to move the device about 2 metres behind point B in order to take a foresight reading from A and a backsight reading from B. Using this configuration, we hoped to measure the degree to which the instruments' performance deviated from the expected.
\end{enumerate}


\subsection*{Part II: Determination of the Stadia Constant (\(K\)) Using the Stadia Method}

\begin{enumerate}
    

\item  \textit{Initial Setup.} In order to create a linear array, we started with a reference point A and spaced out locations B, C, and D; 30 metres apart. Accurate distance measurements were made possible by carefully measuring and marking these distances.

\item  \textit{Instrument Alignment.} By precisely aligning it with the designated point A with the use of a plumb bob, the levelling instrument was set up to guarantee a straight sight line.

\item  \textit{Observation and Recording.} At positions B, C, and D, the readings of the top and bottom stadia wires were recorded with great care. In order to avoid any angular disparities, it was necessary to check the levelling rod's vertical alignment for every measurement.

\item  \textit{Calculations.} The stadia constant (\(K\)) was determined by computing the difference between the top and bottom stadia readings for each location, which was based on the known distance (30 metres) from point A to points B, C, and D.
\end{enumerate}


\section*{Data Analysis}
\addcontentsline{toc}{section}{Data Analysis}


\vspace*{\fill}
\begin{center}
    This area is intentionally left blank.
\end{center}
\vspace{\fill}

\begin{landscape}
\begin{longtable}{ccccc}
\caption{Combined Data for Stadia Method Calculations} \label{table:combined_stadia} \\
\toprule
Distance (m) & Top Hair Reading (cm) & Bottom Hair Reading (cm) & Difference (cm) & Stadia Constant (K) \\
\midrule
\endfirsthead
\caption[]{Combined Data for Stadia Method Calculations} \\
\toprule
Distance (m) & Top Hair Reading (cm) & Bottom Hair Reading (cm) & Difference (cm) & Stadia Constant (K) \\
\midrule
\endhead
\midrule
\multicolumn{5}{r}{Continued on next page} \\
\midrule
\endfoot
\bottomrule
\endlastfoot
20 & 115.500000 & 95.000000 & 20.500000 & 1.025000 \\
30 & 106.500000 & 75.500000 & 31.000000 & 1.033333 \\
\end{longtable}
\end{landscape}

\section*{Results}
\addcontentsline{toc}{section}{Results}


\subsection*{PART I: TWO-PEG TEST}

We painstakingly conducted the Two-Peg Test to see how accurate our levelling device was. This test, albeit basic in nature, revealed important information on the accuracy of the device. Point A's first measurement of 80.2 cm and point B's initial reading of 203.5 cm showed that the technique was understood and that the instrument was reliable when set up appropriately. The repeat measurements provided further confirmation of the instrument's consistency and the correctness of our original results, since they showed little change.


\subsection*{PART II: STADIA METHOD}

In order to get the levelling instrument's stadia constant (K), we used the Stadia Method, which required precise measurements and computations. We measured the stadia constants for each of three places precisely 30 metres from a baseline by carefully noting the top and bottom hair readings. The following constants were determined:

\begin{itemize}
\item[-]  For the distance of 20 meters, a stadia constant (K) of 1.025 was determined.
\item[-] For the distance of 30 meters, a stadia constant (K) of 1.033 was found.
\end{itemize}

A very accurate representation of the instrument's stadia constant is given by the mean K value, which was computed as 1.029. Our results are reliable within a 95\% confidence level, as shown by the confidence interval for this mean K value, which ranges from 0.976 to 1.082. These findings highlight the significance of properly calibrated and operated instruments and the method's efficacy in measuring horizontal distances.

\newpage
\section*{Q \& A}
\addcontentsline{toc}{section}{Q \& A}


\begin{itemize}
\item  \textbf{In the first setup, you were requested to set your level midway of points A and B. What is the main reason?}
   
To reduce systematic mistakes related to the line of sight of the level, it is recommended to position it halfway between points A and B. To ensure that both measurements are accurate, we place the level at an equal distance from both places. This will allow us to compensate for any measurement mistake in the vertical angle. The procedure guarantees an accurate determination of the elevation difference between the two places since any systematic mistake impacting the sight will be equal for both points and cancel out when the elevation difference is calculated. This method successfully removes or at least reduces the influence of any bias in the measuring instruments.


\item \textbf{In the second setup, you were requested to set your level near point B. What is the main reason?}
   
In the second configuration, the level is moved closer to one of the points (here, point B) in order to measure and identify any systematic flaws, such as parallax, that might impact the instrument's accuracy when it is not positioned in the middle of the two locations. This configuration may shed light on the level's performance throughout a range of seeing distances and how the instrument's accuracy changes as a function of distance from the measured location. Surveyors may evaluate and account for any measurement errors caused by the instrument's use at various distances using this error quantification approach.


\item  \textbf{What is the type of error?}
   
In most cases, these configurations will only detect and quantify systematic mistake. Systematic mistakes may be attributed to the instrument's design, setup, or operating procedure and manifest as biases or persistent inaccuracies that happen in the same manner under the same circumstances. A level instrument's sighting mechanism or calibration might be the source of the systematic error that affects the accuracy of elevation readings in the context of the Two-Peg Test. Surveyors may improve the accuracy of their measurements by carrying out the test according to the instructions, which will allow them to detect, measure, and then fix or account for these systematic mistakes.

\end{itemize}


\newpage
\section*{Discussion}
\addcontentsline{toc}{section}{Discussion}



Results from the Stadia Method and Two-Peg Test provide light on the theoretical foundations and practical uses of surveying tools and techniques. Reliability of survey data is ensured by the Two-Peg Test's little fluctuation in repeat readings, which certifies the levelling instrument's accuracy. For more involved or long-distance surveying jobs, even little deviations may add up to big mistakes, therefore accuracy is paramount.

A confidence interval of (0.976, 1.082) was found for the mean stadia constant (K) when the Stadia Method was used. This range indicates that the computed stadia constant is quite precise, which further proves that the device can measure horizontal distances precisely. But a confidence interval shows that there are uncertainties in field measurements, so you have to think about how mistakes may happen and how important it is to be precise while setting up your instruments and taking your measurements.

The little but noticeable variation in stadia constants between 20 and 30 metres demonstrates the complex nature of field surveying, where variables including instrument calibration, ambient circumstances, and operator competence may impact the results of measurements. The estimated mean K value, which is based on these constants, is a reminder of the need for careful surveying techniques and also a proof of how successful the instrument is.

\section*{Conclusion}
\addcontentsline{toc}{section}{Conclusion}

The significance of accuracy and precision in engineering surveying is emphasised in this lab report. By using the Stadia Method and the Two-Peg Test, we proved that the levelling equipment was reliable and that systematic processes are essential for getting accurate surveying results. These activities have given us great practical experience with surveying, and they have shown us how important it is to know the limits and capabilities of our tools.

We now have a far greater respect for the intricacies of field surveying as a result of our work analysing and interpreting the results. In addition to confirming our methodology, the computed stadia constant and confidence interval provide light on the range of variables impacting surveying precision. In order to reduce the occurrence of mistakes and improve the quality of survey data, our experience has highlighted the need of thorough training, ongoing education, and the establishment of industry standards for surveying.

Finally, developing competent surveying techniques relies heavily on the activities described in this paper. By expanding our knowledge and use of surveying concepts, they provide the groundwork for future skill development and have a positive impact on civil engineering as a whole. Going ahead, we must keep doing these kinds of hands-on activities, learning from our mistakes and improving our methods, so that surveying, a vital part of civil engineering, can keep becoming better.



\section*{Appendix}
\addcontentsline{toc}{section}{Appendix}

\begin{lstlisting}[language=Python]
import pandas as pd
import numpy as np
from scipy import stats

data_peg_test = {
    "Measurement": ["Melih Steps", "Initial Staff Readings at Point A", "Initial Staff Readings at Point B"],
    "Description": ["Number of Steps", "Upper (cm)", "Upper (cm)"],
    "Value (cm)": [None, 80.2, 203.5],  # None for steps as it's not directly relevant to calculations here
    "Repeat Measurements (cm)": [None, [80.1, 80.2], [202.5, 202.6]]
}

data_pacing_accuracy = {
    "Measurement": ["Pacing Length", "Pacing Factor"],
    "Value (cm)": [78.466, None],
    "Calculations": [None, 0.7],
    "Average (cm)": [78.466, 0.7]
}

data_stadia_constant = {
    "Distance (m)": [20, 30],
    "Top Hair Reading (cm)": [115.5, 106.5],
    "Bottom Hair Reading (cm)": [95, 75.5],
    "Difference (cm)": [20.5, 31],
    "Stadia Constant (K)": [None, None]  # Placeholder, as it's unclear how to calculate without formula
}

df_peg_test = pd.DataFrame(data_peg_test)
df_pacing_accuracy = pd.DataFrame(data_pacing_accuracy)
df_stadia_constant = pd.DataFrame(data_stadia_constant)

df_stadia_constant['Stadia Constant (K)'] = df_stadia_constant['Difference (cm)'] / df_stadia_constant['Distance (m)']

combined_df = df_stadia_constant

mean_k = combined_df['Stadia Constant (K)'].mean()
print(f"Mean K Value: {mean_k}")

confidence_level = 0.95
degrees_freedom = len(combined_df['Stadia Constant (K)']) - 1
standard_error = stats.sem(combined_df['Stadia Constant (K)'].dropna())

confidence_interval = stats.t.interval(confidence_level, degrees_freedom, mean_k, standard_error)
print(f"Confidence Interval for K Value: {confidence_interval}")

latex_code = combined_df.to_latex(index=False, caption="Combined Data for Stadia Method Calculations", label="table:combined_stadia", longtable=True, column_format="lcccc", header=["Distance (m)", "Top Hair Reading (cm)", "Bottom Hair Reading (cm)", "Difference (cm)", "Stadia Constant (K)"])

print(latex_code)
\end{lstlisting}
\end{document}